\fancyhead[C]{\textit{Solutions}}
\section{\centering Solutions}

\subsection{Solutions to \textit{T/F Questions}}
	\begin{table}[!h]
		\centering
		\begin{tabular}{c@{\hspace{2.5em}}cp{0.75\textwidth}}
			true & false & \\
			\specialrule{\lightrulewidth}{0.25em}{1.25em}
			\bubble & \dbubble & The projection of a vector onto a subspace \(W\) depends on the chosen basis for \(W\).\\[0.75em]
			\dbubble & \bubble & \(\lambda = 1\) is always an eigenvalue of the standard matrix of the transformation that projects vectors in the space
						      onto the subspace \(W\).\\[0.75em]
			\dbubble & \bubble & If \(\vspan \{\vec{a}_1, \vec{a}_2, \vec{a}_3\} = W\), and \(\{\vec{u}_1, \vec{u}_2, \vec{u}_3\}\)
					    is the resulting set of applying the Gram-Schmidt process on the original set of vectors, then 
					    \(\vspan\{\vec{u}_1, \vec{u}_2, \vec{u}_3\} = W\) as well.\\[0.75em]
			\dbubble & \bubble & Every subspace \(W\) of \(\mathbb{R}^n\) has an orthonormal basis.\\[0.75em]
			\bottomrule
		\end{tabular}
	\end{table}


\subsection{Solutions to \textit{Possible/Impossible Questions}}
	\begin{table}[!h]
		\centering
		\begin{tabular}{c@{\hspace{2.5em}}cp{0.75\textwidth}}
			possible & impossible & \\
			\specialrule{\lightrulewidth}{0.25em}{1.25em}
			\dbubble & \bubble & A vector whose projection onto a subspace \(W\) is itself. \\[0.75em]
			\dbubble & \bubble & A non-zero vector whose projection into a subspace is the zero vector.\\[0.75em]
			\bubble & \dbubble & A vector \(\vec{y}\) such that \(\vec{y} - \operatorname{proj}_W \vec{y} \notin W^\perp\).\\[0.75em]
			\bottomrule
		\end{tabular}
	\end{table}


\subsection{Solutions to \textit{Short Answer Questions}}
	\begin{enumerate}
		\item Since \(A\) is an orthogonal matrix, \(\|A\vec{v}\| = \|\vec{v}\| = \sqrt{1^2 + 0^2 + 3^2} = \sqrt{10}\).
		\item We know that for any vector \(\vec{y} \in \mathbb{R}^n\), we can decompose \(\vec{y}\) into the sum of a vector \(\vec{\hat{y}} \in W\) and
			a vector \(\vec{z} \in W^\perp\). That is, we can always decompose any vector \(\vec{y}\) as \[\vec{y} = \vec{\hat{y}} + \vec{z}.\] We know that
			\(\vec{\hat{y}} = \proj_W\vec{y}\), and we're given what the projected vector is. All we have to do now is to find an appropriate vector 
			\(\vec{z} \ne \vec{0}\). The vector \(\vec{z}\) must satisfy the system of linear equations 
			\[\left\{\begin{aligned}
				\vec{z} &\cdot \vec{u}_1 = 0 \\
				\vec{z} &\cdot \vec{u}_2 = 0
			\end{aligned}\right. \iff \left\{\begin{aligned}z_1 + z_2 + z_3 &= 0  \\ -z_1 + 0z_2 + z_3 &= 0\end{aligned}\right.\]
			We find that any \(\vec{z} = t\begin{bmatrix}-1 \\ -2 \\ 1\end{bmatrix}\) works, for any \(t \in \mathbb{R}\). So a valid answer is 
			\[\vec{y} = (2\vec{u}_1 + 3\vec{u}_2) + \begin{bmatrix}-1 \\ -2 \\ 1\end{bmatrix},\] and in general we can multiply the rightmost vector by 
			any real scalar to get a new vector that satisfies this problem.

			To answer the ``challenge" problem, we just have to make sure that \(\|\vec{z}\| = 3\). So, we need that 
			\[\sqrt{t^2 + 4t^2 + t^2} = 3 \iff \sqrt{6t^2} = 3 \iff \lvert t \rvert = \frac{3}{\sqrt{6}}.\] So both 
			\[\vec{y} = (2\vec{u}_1 + 3\vec{u}_2) + \frac{3}{\sqrt{6}}\begin{bmatrix}-1 \\ -2 \\ 1\end{bmatrix}, \quad \vec{y} = (2\vec{u}_1 + 3\vec{u}_2) - \frac{3}{\sqrt{6}}\begin{bmatrix}-1 \\ -2 \\ 1\end{bmatrix}\] satisfy the constraints.

		\item This is just \(\proj_P \vec{v} = \begin{bmatrix}3 \\ 5 \\ 3\end{bmatrix}\). Don't forget to find an orthogonal basis for \(P\)!
		\item This new basis is orthogonal, so the typical formula for orthogonal projection could be used! You'll find that you get the same answer :)
	\end{enumerate}

